% vim: set textwidth=78 autoindent:

\section{Usar complementos del núcleo de QGIS}\label{sec:core_plugins}\index{plugins!core}

% when the revision of a section has been finalized, 
% comment out the following line:
%\updatedisclaimer

QGIS contiene actualmente 19 complementos del núcleo que se pueden cargar usando el Administrador de complementos.
La tabla \ref{tab:core_plugins} lista cada uno de los complementos del núcleo junto con una
descripción de su propósito y el icono de la barra de herramientas.\footnote{El complemento
de exportación a MapServer y el Instalador de complementos son complementos externos de Python,
pero forman parte de las fuentes de QGIS y se cargan automáticamente y se pueden seleccionar 
desde el Administrador de complementos de QGIS.}

% minipage is needed to appear the footnote under the table
% SH
\begin{minipage}{\textwidth}
\begin{table}[H]
\centering
\caption{Complementos del núcleo de QGIS}\label{tab:core_plugins}\medskip
\small
 \begin{tabular}{|l|l|p{4in}|}
\hline \textbf{Icono} & \textbf{Complemento} & \textbf{Descripción}\\
\hline
\includegraphics[width=0.6cm]{delimited_text}
 & Añadir capa de texto delimitado \index{plugins!delimited text} & Carga y muestra archivos de texto delimitado que contengan coordenadas X e Y\\
\hline
\includegraphics[width=0.6cm]{coordinate_capture}
 & Captura de coordenadas \index{plugins!coordinate capture}& Captura las coordenadas del ratón en diferentes SRC\\
\hline 
\includegraphics[width=0.6cm]{copyright_label}
 & Etiqueta de Copyright \index{plugins!copyright}& Dibuja una etiqueta de copyright con información\\
\hline 
\includegraphics[width=0.6cm]{dxf2shp_converter}
 & Conversor DXFaSHP \index{plugins!DXF2Shape}& Convierte archivos DXF a formato SHP\\
\hline
\includegraphics[width=0.6cm]{gps_importer}
 & Herramientas de GPS \index{plugins!gps}& Herramientas para cargar e importar datos de GPS\\
\hline
\includegraphics[width=0.6cm]{grass}
 & GRASS \index{plugin!grass toolbox} & Activa la poderosa Caja de Herramientas de GRASS\\
\hline
\includegraphics[width=0.6cm]{georeferencer}
 & Georreferenciador \index{plugin!georeferencer} & Añade información de proyección a archivos ráster\\
\hline
\includegraphics[width=0.6cm]{grid_maker}
 & Creador de cuadrícula \index{plugins!graticule}& Crea una cuadrícula latitud/longitud y la guarda como archivo shape\\
\hline
\includegraphics[width=0.6cm]{interpolation}
& Complemento de interpolación \index{plugins!Interpolation}& Interpolación en base a los vértices de una capa vectorial\\
\hline
\includegraphics[width=0.6cm]{mapserver_export}
& Complemento de exportación a MapServer \index{plugins!MapServer Export}& Exporta un proyecto guardado de QGIS a un archivo de mapa de MapServer \\
\hline
\includegraphics[width=0.6cm]{north_arrow}
& Flecha de Norte \index{plugins!north arrow}& Muestra una flecha de Norte superpuesta sobre el mapa\\
\hline
\includegraphics[width=0.6cm]{ogr_converter}
 & Conversor de capas OGR \index{plugins!OGR converter} & Traduce capas vectoriales entre formatos admitidos por OGR\\
\hline
\includegraphics[width=0.6cm]{plugin_installer}
 & Instalador de complementos \index{plugins!Plugin Installer} & Descarga e instala complementos de python de QGIS\\
\hline
\includegraphics[width=0.6cm]{spiticon}
 & SPIT \index{plugins!spit}& Herramienta de importación de archivos shape a PostgreSQL/PostGIS \\
\hline
\includegraphics[width=0.6cm]{quick_print}
 & Impresión rápida \index{plugins!quick print}& Imprimir un mapa rápidamente con un mínimo esfuerzo \\
\hline
\includegraphics[width=0.6cm]{scale_bar}
 & Barra de escala \index{plugins!scalebar}& Dibuja una barra de escala\\
\hline
\includegraphics[width=0.6cm]{mIconAddWfsLayer}
 & WFS & Carga y muestra capas WFS \\
\hline
\end{tabular}
\end{table}
\end{minipage}

\normalsize

\begin{Tip}\caption{\textsc{Configuración de complementos guardada en proyecto}}\index{plugins settings}
\qgistip{Cuando guarda un proyecto .qgs, cualquier cambio que haya hecho en los complementos Fecha de Norte, Barra de escala y Etiqueta de Copyright se guardará en el proyecto y se recuperará la próxima vez que cargue el proyecto.}
\end{Tip}
